% ==============================================================================
% SECCIÓN 05.01: MUESTREO ALEATORIO SIMPLE, MEDIA Y VARIANZA MUESTRAL INSESGADA
% ==============================================================================

\subsection*{Enunciados de los problemas --- Sección 05.01}

\subsubsection*{Nivel Fundamental}

\begin{problema}
	\label{prob:5.1.1}
	Una muestra aleatoria simple de tamaño $n=5$ produce las observaciones $\{8, 12, 10, 14, 6\}$.
	\begin{enumerate}
		\item Calcule la media muestral $\bar{X}$.
		\item Calcule la varianza muestral insesgada $S^{2}$.
	\end{enumerate}
\end{problema}

\begin{problema}
	\label{prob:5.1.2}
	Los puntajes de un examen estandarizado tienen media poblacional $\mu = 500$ y desviación estándar $\sigma = 40$. Se toma una muestra aleatoria simple de $n = 64$ estudiantes.
	\begin{enumerate}
		\item Calcule $E(\bar{X})$.
		\item Calcule $\Var(\bar{X})$ y el error estándar $\text{DE}(\bar{X})$.
	\end{enumerate}
\end{problema}

\begin{problema}
	\label{prob:5.1.3}
	Clasifique cada una de las siguientes cantidades como \emph{estadístico} o \emph{parámetro}, justificando su respuesta: (a) la media poblacional $\mu$; (b) la media muestral $\bar{x}$ de una encuesta; (c) la proporción muestral $\hat{p}$ de una elección; (d) la varianza poblacional $\sigma^{2}$.
\end{problema}

\subsubsection*{Nivel Operativo}

\begin{problema}
	\label{prob:5.1.4}
	Sean $X_{1}, \ldots, X_{n}$ variables aleatorias independientes con la misma varianza $\sigma^{2}$. Partiendo de la propiedad $\Var(\sum_i X_i) = \sum_i \Var(X_i)$ para variables independientes, derive la fórmula $\Var(\bar{X}) = \sigma^{2}/n$.
\end{problema}

\begin{problema}
	\label{prob:5.1.5}
	Un proceso de medición tiene desviación estándar poblacional conocida $\sigma = 15$. Calcule el error estándar $\text{DE}(\bar{X}) = \sigma/\sqrt{n}$ para tamaños de muestra $n = 25$, $n = 100$ y $n = 400$, y verifique que cuadruplicar $n$ reduce el error estándar a la mitad.
\end{problema}

\begin{problema}
	\label{prob:5.1.6}
	Para los datos $\{5, 7, 9, 11, 13\}$ ($n=5$):
	\begin{enumerate}
		\item Calcule la varianza muestral insesgada $S^{2} = \frac{1}{n-1}\sum(X_i-\bar{X})^{2}$.
		\item Calcule el estimador ``natural'' (sesgado) $\hat{\sigma}^{2} = \frac{1}{n}\sum(X_i-\bar{X})^{2}$.
		\item Verifique numéricamente la relación $\hat{\sigma}^{2} = \frac{n-1}{n}S^{2}$.
	\end{enumerate}
\end{problema}

\subsubsection*{Nivel Analítico}

\begin{problema}
	\label{prob:5.1.7}
	Demuestre que $E(S^{2}) = \sigma^{2}$ para $X_{1}, \ldots, X_{n}$ iid con media $\mu$ y varianza $\sigma^{2}$. Justifique explícitamente por qué $\sum_{i=1}^{n}(X_i-\mu) = n(\bar{X}-\mu)$ y cómo se usa esta identidad en la demostración.
\end{problema}

\begin{problema}
	\label{prob:5.1.8}
	Demuestre formalmente que $\Var(\bar{X}) = \sigma^{2}/n$ usando la independencia de $X_{1}, \ldots, X_{n}$, y explique por qué este resultado, junto con la Ley de los Grandes Números, implica que $\bar{X}$ es un \emph{estimador consistente} de $\mu$ (es decir, que $\Var(\bar{X}) \to 0$ cuando $n \to \infty$).
\end{problema}

\subsubsection*{Nivel Desafiante}

\begin{problema}
	\label{prob:5.1.9}
	Demuestre algebraicamente la fórmula de cálculo directo (``fórmula abreviada'') de la varianza muestral: $S^{2} = \frac{1}{n-1}\left[\sum_{i=1}^n X_i^{2} - n\bar{X}^{2}\right]$, partiendo de la definición $S^{2} = \frac{1}{n-1}\sum(X_i-\bar{X})^{2}$. Aplique ambas fórmulas a los datos $\{2, 4, 6, 8, 10\}$ y verifique que producen el mismo resultado.
\end{problema}

\begin{problema}
	\label{prob:5.1.10}
	Una población finita de $N = 200$ cuentas bancarias tiene varianza $\sigma^{2} = 900$. Se extrae una muestra aleatoria simple de $n = 20$ cuentas \emph{sin reemplazo}.
	\begin{enumerate}
		\item Calcule $\Var(\bar{X})$ usando la fórmula ingenua de población infinita $\sigma^{2}/n$.
		\item Calcule $\Var(\bar{X})$ usando el factor de corrección por población finita (FPC): $\Var(\bar{X}) = \frac{\sigma^{2}}{n}\cdot\frac{N-n}{N-1}$.
		\item Explique por qué la FPC siempre reduce la varianza cuando se muestrea sin reemplazo de una población finita.
	\end{enumerate}
\end{problema}

\subsection*{Sugerencias de los problemas --- Sección 05.01}

\begin{sugerencia}
	\label{sug:5.1.1}
	Para el Problema \ref{prob:5.1.1}, $\bar{X} = 50/5 = 10$. Las desviaciones al cuadrado son $4, 4, 0, 16, 16$, con suma $40$; divida entre $n-1=4$.
\end{sugerencia}

\begin{sugerencia}
	\label{sug:5.1.2}
	Para el Problema \ref{prob:5.1.2}, use directamente $E(\bar{X})=\mu$ y $\Var(\bar{X})=\sigma^{2}/n$ con $\sigma^{2}=1600$ y $n=64$.
\end{sugerencia}

\begin{sugerencia}
	\label{sug:5.1.3}
	Para el Problema \ref{prob:5.1.3}, un parámetro describe a la población (fijo, generalmente desconocido); un estadístico se calcula a partir de una muestra observada (variable aleatoria, conocido una vez tomada la muestra).
\end{sugerencia}

\begin{sugerencia}
	\label{sug:5.1.4}
	Para el Problema \ref{prob:5.1.4}, escriba $\bar{X}=\frac{1}{n}\sum X_i$ y use $\Var(aY)=a^{2}\Var(Y)$ junto con la aditividad de la varianza bajo independencia.
\end{sugerencia}

\begin{sugerencia}
	\label{sug:5.1.5}
	Para el Problema \ref{prob:5.1.5}, $\text{DE}(\bar{X})=15/\sqrt{n}$: para $n=25$, $15/5=3$; para $n=100$, $15/10=1.5$; para $n=400$, $15/20=0.75$.
\end{sugerencia}

\begin{sugerencia}
	\label{sug:5.1.6}
	Para el Problema \ref{prob:5.1.6}, $\bar{X}=9$; las desviaciones al cuadrado son $16,4,0,4,16$, con suma $40$. Divida entre $4$ para $S^{2}$ y entre $5$ para $\hat{\sigma}^{2}$.
\end{sugerencia}

\begin{sugerencia}
	\label{sug:5.1.7}
	Para el Problema \ref{prob:5.1.7}, siga la demostración del Teorema de insesgadez de $S^2$: escriba $(X_i-\bar X)=(X_i-\mu)-(\bar X-\mu)$, expanda el cuadrado, sume sobre $i$ y use $\sum(X_i-\mu)=n(\bar X-\mu)$ (que se sigue de la definición de $\bar X$) para simplificar el término cruzado.
\end{sugerencia}

\begin{sugerencia}
	\label{sug:5.1.8}
	Para el Problema \ref{prob:5.1.8}, $\Var(\bar{X})=\Var\left(\frac{1}{n}\sum X_i\right)=\frac{1}{n^{2}}\sum\Var(X_i)=\frac{n\sigma^{2}}{n^{2}}=\sigma^{2}/n \to 0$ cuando $n\to\infty$.
\end{sugerencia}

\begin{sugerencia}
	\label{sug:5.1.9}
	Para el Problema \ref{prob:5.1.9}, expanda $(X_i-\bar X)^2 = X_i^2 - 2X_i\bar X + \bar X^2$, sume sobre $i$, y use $\sum X_i = n\bar X$ para simplificar $-2\bar X\sum X_i + n\bar X^2$ a $-n\bar X^2$.
\end{sugerencia}

\begin{sugerencia}
	\label{sug:5.1.10}
	Para el Problema \ref{prob:5.1.10}, el factor FPC $(N-n)/(N-1) = 180/199 \approx 0.9045$ es siempre menor que $1$ cuando $n>1$ y $N$ es finito, por lo que siempre reduce la varianza ingenua.
\end{sugerencia}

\subsection*{Soluciones de los problemas --- Sección 05.01}

\begin{solucion}
	\label{sol:5.1.1}
	Para el Problema \ref{prob:5.1.1} con datos $\{8, 12, 10, 14, 6\}$:
	\begin{enumerate}
		\item $\bar{X} = \dfrac{8+12+10+14+6}{5} = \dfrac{50}{5} = 10$.
		\item Las desviaciones al cuadrado son $(8-10)^{2}=4$, $(12-10)^{2}=4$, $(10-10)^{2}=0$, $(14-10)^{2}=16$, $(6-10)^{2}=16$, con suma $40$. Por lo tanto, $S^{2} = 40/(5-1) = 10$. \quad \qedhere
	\end{enumerate}
\end{solucion}

\begin{solucion}
	\label{sol:5.1.2}
	Para el Problema \ref{prob:5.1.2} con $\mu=500$, $\sigma=40$, $n=64$:
	\begin{enumerate}
		\item $E(\bar{X}) = \mu = 500$.
		\item $\Var(\bar{X}) = \sigma^{2}/n = 1600/64 = 25$. $\text{DE}(\bar{X}) = \sqrt{25} = 5$. \quad \qedhere
	\end{enumerate}
\end{solucion}

\begin{solucion}
	\label{sol:5.1.3}
	Para el Problema \ref{prob:5.1.3}: (a) $\mu$ es un \textbf{parámetro} (describe a toda la población, es fijo y generalmente desconocido). (b) $\bar{x}$ es un \textbf{estadístico} (se calcula a partir de la muestra observada). (c) $\hat{p}$ es un \textbf{estadístico} (proporción calculada de datos muestrales). (d) $\sigma^{2}$ es un \textbf{parámetro} (describe la dispersión de toda la población). \quad \qedhere
\end{solucion}

\begin{solucion}
	\label{sol:5.1.4}
	Para el Problema \ref{prob:5.1.4}:
	\begin{align*}
		\Var(\bar{X}) = \Var\left(\frac{1}{n}\sum_{i=1}^{n} X_i\right) = \frac{1}{n^{2}}\Var\left(\sum_{i=1}^{n} X_i\right)
		= \frac{1}{n^{2}}\sum_{i=1}^{n}\Var(X_i) = \frac{1}{n^{2}}\cdot n\sigma^{2} = \frac{\sigma^{2}}{n}. \quad \qedhere
	\end{align*}
\end{solucion}

\begin{solucion}
	\label{sol:5.1.5}
	Para el Problema \ref{prob:5.1.5} con $\sigma=15$:
	\begin{itemize}
		\item $n=25$: $\text{DE}(\bar{X}) = 15/\sqrt{25} = 15/5 = 3$.
		\item $n=100$: $\text{DE}(\bar{X}) = 15/\sqrt{100} = 15/10 = 1.5$.
		\item $n=400$: $\text{DE}(\bar{X}) = 15/\sqrt{400} = 15/20 = 0.75$.
	\end{itemize}
	Cuadruplicar $n$ ($25\to100$, $100\to400$) reduce el error estándar exactamente a la mitad ($3\to1.5$, $1.5\to0.75$), consistente con la dependencia $1/\sqrt{n}$. \quad \qedhere
\end{solucion}

\begin{solucion}
	\label{sol:5.1.6}
	Para el Problema \ref{prob:5.1.6} con datos $\{5,7,9,11,13\}$, $\bar{X}=9$, y desviaciones al cuadrado $16,4,0,4,16$ (suma $40$):
	\begin{enumerate}
		\item $S^{2} = 40/(5-1) = 10$.
		\item $\hat{\sigma}^{2} = 40/5 = 8$.
		\item Verificación: $\frac{n-1}{n}S^{2} = \frac{4}{5}\cdot 10 = 8 = \hat{\sigma}^{2}$. \quad \qedhere
	\end{enumerate}
\end{solucion}

\begin{solucion}
	\label{sol:5.1.7}
	Para el Problema \ref{prob:5.1.7}: como $\bar{X}=\frac{1}{n}\sum X_i$, se sigue que $\sum_{i=1}^n(X_i-\mu) = \sum X_i - n\mu = n\bar{X}-n\mu = n(\bar{X}-\mu)$. Usando esta identidad,
	\begin{align*}
		\sum_{i=1}^n(X_i-\bar X)^2 &= \sum_{i=1}^n\left[(X_i-\mu)-(\bar X-\mu)\right]^2 \\
		&= \sum_{i=1}^n(X_i-\mu)^2 - 2(\bar X-\mu)\underbrace{\sum_{i=1}^n(X_i-\mu)}_{=n(\bar X-\mu)} + n(\bar X-\mu)^2 \\
		&= \sum_{i=1}^n(X_i-\mu)^2 - n(\bar X-\mu)^2.
	\end{align*}
	Tomando esperanza, $E[(X_i-\mu)^2]=\sigma^2$ y $E[(\bar X-\mu)^2]=\Var(\bar X)=\sigma^2/n$, así que $E\left[\sum(X_i-\bar X)^2\right] = n\sigma^2 - n\cdot\sigma^2/n = (n-1)\sigma^2$. Dividiendo entre $n-1$: $E(S^2)=\sigma^2$. \quad \qedhere
\end{solucion}

\begin{solucion}
	\label{sol:5.1.8}
	Para el Problema \ref{prob:5.1.8}: por independencia de $X_1,\ldots,X_n$,
	\begin{align*}
		\Var(\bar X) = \frac{1}{n^2}\sum_{i=1}^n \Var(X_i) = \frac{n\sigma^2}{n^2} = \frac{\sigma^2}{n}.
	\end{align*}
	Como $\sigma^2$ es finita y fija, $\Var(\bar X) = \sigma^2/n \to 0$ cuando $n\to\infty$. Junto con $E(\bar X)=\mu$ para todo $n$ (insesgadez), esto implica que $\bar X$ converge en probabilidad a $\mu$ (por la desigualdad de Chebyshev), es decir, $\bar X$ es un estimador \emph{consistente} de $\mu$. \quad \qedhere
\end{solucion}

\begin{solucion}
	\label{sol:5.1.9}
	Para el Problema \ref{prob:5.1.9}: expandiendo el cuadrado,
	\begin{align*}
		\sum_{i=1}^n(X_i-\bar X)^2 = \sum_{i=1}^n\left(X_i^2 - 2X_i\bar X + \bar X^2\right)
		= \sum_{i=1}^n X_i^2 - 2\bar X\sum_{i=1}^n X_i + n\bar X^2.
	\end{align*}
	Como $\sum X_i = n\bar X$, el término cruzado es $-2\bar X\cdot n\bar X = -2n\bar X^2$, así que la expresión se reduce a $\sum X_i^2 - 2n\bar X^2 + n\bar X^2 = \sum X_i^2 - n\bar X^2$. Por lo tanto $S^2 = \frac{1}{n-1}\left[\sum X_i^2 - n\bar X^2\right]$.

	Para $\{2,4,6,8,10\}$: $\bar X = 30/5=6$. Fórmula directa: desviaciones al cuadrado $16,4,0,4,16$, suma $40$, $S^2=40/4=10$. Fórmula abreviada: $\sum X_i^2 = 4+16+36+64+100=220$, $n\bar X^2 = 5\cdot 36=180$, $S^2=(220-180)/4=40/4=10$. Ambas coinciden. \quad \qedhere
\end{solucion}

\begin{solucion}
	\label{sol:5.1.10}
	Para el Problema \ref{prob:5.1.10} con $N=200$, $\sigma^2=900$, $n=20$:
	\begin{enumerate}
		\item Fórmula ingenua: $\Var(\bar X) = \sigma^2/n = 900/20 = 45$.
		\item Con FPC: $\Var(\bar X) = \dfrac{\sigma^2}{n}\cdot\dfrac{N-n}{N-1} = 45\cdot\dfrac{180}{199} \approx 45\cdot 0.9045 \approx 40.70$.
		\item El factor $(N-n)/(N-1)$ es estrictamente menor que $1$ siempre que $n>1$ y $N$ sea finito (pues $N-n<N-1$ cuando $n>1$), por lo que la FPC siempre reduce la varianza estimada respecto al caso de población infinita; cuando $N\to\infty$ (o $n/N\to 0$), el factor tiende a $1$ y ambas fórmulas coinciden. \quad \qedhere
	\end{enumerate}
\end{solucion}
